\setauthor{\secondauthor}

\section{System Requirements and Constraints}

\subsection{Real-Time Processing Requirements}

\subsubsection{Latency Constraints}
The entire process from capture, through preprocessing, motion detection, blob extraction, tip detection, and scoring, needs to be completed within the time interval between two successive frames. Since the frame rate is 30 frames per second, the time constraint is 33 milliseconds, while at 15 frames per second, the time constraint is 66 milliseconds \cite{Szeliski2022}. It was decided that the end-to-end latency from the time the dart hits until the time the score appears on the screen needs to be less than 100 milliseconds. This was determined through informal testing, where it was observed that such high delays are noticeable during the course of play. Non-critical processing such as the creation of overlays and match logging is done after the commitment of the detection result and is not subject to the above constraints.

\subsubsection{Pipeline Responsiveness}
The responsiveness of the pipeline is maintained by making the more expensive operations dependent on the results obtained from the less expensive operations. Capture of the frame and the differencing operation are performed unconditionally on every frame. Subpixel tip detection, perspective correction, and sector assignment are performed on every frame, provided the differencing operation indicates motion greater than the specified threshold.


\subsection{Hardware Constraints}

\subsubsection{Single-Camera Assumption}
The system assumes the availability of a single camera.
Restricting the problem domain to a single-camera setup was a conscious decision. This eliminates synchronization problems, keeps the calibration process manageable for the end user, and minimizes hardware costs to a single USB camera module. In practice, spatial reasoning is limited to a single image plane. A hand occluding a dart or a previously lodged dart cannot be resolved geometrically. However, this is acceptable for the problem domain because the dart tip is typically visible from a frontal overhead position at the time of impact.

\subsubsection{Static Mount Requirements}
This detection approach is predicated on a static camera position between the time of calibration and use. This is necessary for the background model, the stored ArUco pose, and the pixel-to-board homography. Any physical movement of the camera, even slight, invalidates the stored calibration and causes a scoring problem due to a systematic error. A startup routine is implemented to ensure the presence of a valid stored calibration. If this is missing or marked as invalid, the system prompts a recalibration routine.

\section{Programming Language and Runtime Environment}

\subsection{Python as Core Development Language}

\subsubsection{Computer Vision Suitability}
Python has been selected due to the maturity of the OpenCV bindings. Every function in the performance-critical path, i.e., encompassing the camera grab, color conversion, morphology operations, contour detection, ArUco detection, and \texttt{solvePnP} execution, is performed within OpenCV's compiled C++ backend. Python is used to orchestrate this process. Hence, only function call overhead is incurred in the Python interpreter. This is consistent with the cost model we would expect for this type of real-time processing pipeline.


\subsubsection{Scientific Ecosystem}
This is because the underlying data structure of NumPy arrays matches the internal memory structure of OpenCV, allowing for efficient passing of frames between libraries without copying data. Additionally, the scipy library offered the necessary geometric primitives for implementing the mapping using polar coordinates during the early stages of prototyping. Matplotlib also helped in the inspection of intermediate results during the development stage, such as difference frames, binary masks, and contour maps, without needing a separate debugging tool. The fact that the above libraries are available as drop-in modules greatly reduced the development cycles during parameter tuning.

\subsection{Runtime Characteristics}

\subsubsection{Interpreter-Based Execution}
The interpreter imposes a certain overhead on the execution of the Python bytecode in a loop over the pixel data. However, this was not a performance constraint, as no stage of the pipeline needs to loop over the pixel data at the Python level. All pixel array operations are vectorized using NumPy or OpenCV \cite{Harris2020}. There is one loop at the Python level, which iterates over a small list of filtered contour candidates. This has negligible performance impact, as it needs to process at most a few objects per frame.


\subsubsection{Performance Considerations}
The first optimization step for the project involved the implementation of region of interest cropping. This involved limiting the processing of the image to the region of interest, which contains the dartboard. This resulted in a reduction of between 60-70 percent of the total number of pixels processed compared to the whole frame \cite{Viola2001}. The output buffers are also pre-filled, which are then passed as arguments for OpenCV's in-place processing, thus eliminating the need for heap allocation. In the event that the early exit condition is met, with no motion detected, the processing pipeline takes a short time, less than 10 milliseconds, thus giving ample time for the UI thread.

\begin{center}
\captionsetup{type=table,hypcap=false}
	\centering
	\caption{Approximate per-stage runtime budget interpretation for practical operation.}
	\label{tab:runtime_budget}
	\begin{tabularx}{\linewidth}{@{}>{\RaggedRight\arraybackslash}X >{\centering\arraybackslash}p{3.2cm} >{\RaggedRight\arraybackslash}X@{}}
		\toprule
		\shortstack[l]{\textbf{Pipeline}\\\textbf{Stage}} & \shortstack{\textbf{Typical}\\\textbf{Cost Class}} & \shortstack[l]{\textbf{Optimization}\\\textbf{Used}} \\
		\midrule
		Frame acquisition + grayscale conversion & Low & Direct OpenCV capture path \\
		Differencing + thresholding + morphology & Medium & ROI restriction and in-place operations \\
		Contour extraction + tip estimation & Medium & Early exit when no motion is present \\
		Score mapping + event emission & Low & Lightweight geometric classification \\
		Overlay/debug composition & Low to Medium & Performed after detection commitment \\
		\bottomrule
	\end{tabularx}
\end{center}

\section{Computer Vision Framework}

\subsection{OpenCV}

\subsubsection{Image Capture and Preprocessing}
The images are captured using the OpenCV function \texttt{cv2.VideoCapture}. After that, the images are converted to grayscale, which is necessary for the detection. However, sometimes it is also necessary to keep the color information, for example, to filter the specular highlights on the metallic ring. In this case, the image is converted to the HSV space. Lens distortion is compensated using the intrinsic parameters obtained during the chessboard calibration. This is necessary for the outer scoring rings, where otherwise the apparent position of the darts on the board would be off by several millimeters due to the barrel distortion.

\subsubsection{Filtering and Thresholding}
A Gaussian blur is applied before the background subtraction step to remove sensor noise which could cause unwanted contour formation. Otsu's method is then used to set the binary threshold after background subtraction, which eliminates the need for retuning the fixed threshold value in response to changing ambient light conditions \cite{Otsu1979}. Morphological opening is then used to remove isolated single-pixel noise, which is followed by a closing function to fill in small gaps in dart-sized objects before contour extraction. As a last resort for dealing with environments featuring significant illumination gradients across the board, adaptive thresholding is also supported, though at the expense of a small increase in processing time.

\subsection{Supporting Libraries}

\subsubsection{NumPy for Numerical Computation}
All coordinate data computed during the OpenCV pipeline, including contour points, homography, and marker corners, is represented as NumPy arrays and remains this way throughout the entire process. The computation of pixel coordinates from board coordinates is accomplished through a matrix multiplication followed by an \texttt{arctan2} function, which is performed over the entire set of candidate points simultaneously. This maintains the high speed of the coordinate computation step even when multiple candidates are present.

\subsubsection{Geometry Utility Libraries}
A utility library has been developed for handling the conversion between polar coordinates and sector coordinates, including identifying which numbered segment and which type of ring (single, double, triple, or bull) a given board coordinate represents. It uses a lookup table of angular and radial thresholds to represent the standard dartboard layout and exposes a single function which accepts a coordinate in board space and returns the corresponding score value. This separation of code from the detection routine also allowed for the independent unit testing of the dartboard scoring rules without the requirement for image resources.

\section{Calibration Approaches}

\subsection{Marker-Based Calibration}

\subsubsection{ArUco Marker Detection}
Four ArUco markers are placed at known positions around the board frame \cite{Garrido-Jurado2014}. The corner points are identified in the camera view through \texttt{cv2.aruco.detectMarkers} with the ArUco dictionary \texttt{DICT\_4X4\_50}. The markers are uniquely identified with different IDs. Thus, even in the event that one marker is occluded, the corner points are identified unambiguously. The corner points are then passed to the next stage for estimation.


\subsubsection{Pose Estimation}
The corner points are then used to obtain the camera pose with respect to the board through \texttt{cv2.solvePnP}. A homography matrix is then obtained that maps any point in the camera view to the corresponding point in the front view \cite{Hartley2004}. The homography matrix is stored for the entire duration as the camera is not moving. 

\subsection{Manual Calibration}

\subsubsection{User-Defined Ring and Sector Placement}
The user may choose not to use the ArUco markers for calibration. Instead, the user is prompted to click the center of the bullseye in the live view window. Then the user is prompted to click any point on the outer double ring. The two points are used to calculate the center and radius of the dartboard in pixels. All the radial thresholds are then scaled from the derived radius according to the proportions in the conventional dartboard.
After the rings are set up, the user selects one of the known sector boundaries. The known sector boundary between segments 20 and 1 located at the top of the board was used as the standard to measure the angular offset. The 20 sector boundaries are then evenly distributed with 18-degree increments of the reference angle. This is necessary since the board's rotational orientation with regard to the camera cannot be determined beforehand and will change depending on how the board is mounted.

\subsection{Comparison of Calibration Methods}

\subsubsection{Setup Complexity}
ArUco calibration requires the user to place four markers and attach them to the board before first-time usage. Once the markers are attached, the process runs automatically every time the board starts without further user intervention. Manual calibration requires no pre-installation but requires the user to interact with the program several times each time the board starts or the camera is moved. For a permanent installation, ArUco would be much simpler to set up; for an ad hoc setting, manual calibration would be preferable to the requirement of attaching four markers to the board.

\subsubsection{Robustness in Real Environments}
Based on the experiments performed, it was found that the ArUco method of calibration was more reliable and consistent with regard to the assignment of each sector between sessions since it takes into consideration the homography to correct for perspective distortion. The manual method would gradually accumulate error depending on the accuracy of the clicks and would not compensate for perspective problems. Although under controlled conditions the accuracy of the ArUco and manual method of calculating the homography would be nearly equivalent with regard to the accuracy of the points scored by each method, under less than optimal camera angles with regard to the board, the homography method would clearly produce fewer errors with regard to the correct classification of adjacent segments along the board boundaries.

\vspace{\baselineskip}

\begin{center}
\captionsetup{type=table,hypcap=false}
	\centering
	\caption{Operational comparison of calibration strategies in this project.}
	\label{tab:calibration_strategy_comparison}
	\begin{tabularx}{\linewidth}{@{}X X X@{}}
		\toprule
		\textbf{Criterion} & \textbf{Manual Calibration} & \textbf{ArUco Calibration} \\
		\midrule
		Initial setup effort & Low hardware, higher interactive effort & Marker placement required once \\
		Repeatability across sessions & Moderate, click-dependent & High, marker pose-dependent \\
		Perspective compensation quality & Limited in oblique views & Strong through homography estimation \\
		Best fit scenario & Ad hoc temporary setups & Fixed or semi-permanent installations \\
		\bottomrule
	\end{tabularx}
\end{center}

\vspace{\baselineskip}

\section{Motion Detection Techniques}

\subsection{Frame Differencing}

\subsubsection{Background Selection}
The background model is pre-initialized using a median value computed over a short frame buffer taken before play begins, during which time no dart is on the board. A median rather than a mean value is used to resist transient effects of passing shadows or exposure changes \cite{Stauffer1999}. During play, the background model is not updated since the dart should be visible as a persistent foreground object. Background updating is invoked only when the user indicates that the darts have been removed from the board.

\subsubsection{Temporal Stability}
A single-frame difference is noisy by nature. A dart landing in a single frame and staying there in subsequent frames would result in a high difference in the landing frame and almost zero differences in subsequent frames. To avoid false alarms due to fluorescent lighting flicker and board surface vibrations, the system requires the blob to be present in at least two consecutive difference frames before it is considered a hit.


\subsection{Threshold-Based Segmentation}

\subsubsection{Binary Mask Generation}
Otsu's thresholding is used to convert the absolute difference between the current frame and the background model to a binary mask. This is quite applicable in this case, considering that the difference image is mostly dark with a small bright blob corresponding to the dart. This is a binary image problem, and Otsu's method is applicable in this case. This mask is the only input to the contour extraction process.

\subsubsection{Sensitivity to Noise}
There are two sources of noise in this problem. One is the specular reflection from the metallic ring between the double and treble regions. The other is pixel noise at low levels. Sensor noise is taken care of by the Gaussian blur applied before the differencing process. Specular reflection is taken care of by the morphological opening process that removes objects with sizes below a specified minimum area. This is set empirically based on the false positives encountered on the development board.


\subsection{Blob Detection}

\subsubsection{Contour Extraction}
The function \texttt{cv2.findContours} is called with the binary mask and the flag
\texttt{RETR\_EXTERNAL} to retrieve only the outermost contours while ignoring any
nested ones. This is sufficient for the dart tip projection because it is a compact
area with no inner regions. The returned value is a list of contour points that
describe the outline of each region in the image.


\subsubsection{Filtering of Irrelevant Regions}
The regions are filtered by four checks before any of the tip coordinates are computed. The first check is that of area. Regions that are either below the empirically determined minimum or exceed the maximum possible cross-section area are eliminated. The second check involves the aspect ratio. Highly elongated regions are unlikely to correspond to a dart tip when viewed from above. The third check involves the shape solidity, which is given by the ratio of the contour area to the area of the convex hull. Irregular shapes resulting from reflections or cables have lower solidity values. The fourth check involves the spatial location of the region. Regions whose centroids are outside the known region of the board are eliminated. The number of regions that remain at this point per frame is usually either zero or one, occasionally two, before the tip localization.

\begin{center}
\captionsetup{type=table,hypcap=false}
	\centering
	\caption{Key detector parameters and operational role.}
	\label{tab:detector_parameters}
	\begin{tabularx}{\linewidth}{@{}X X X@{}}
		\toprule
		\textbf{Parameter} & \textbf{Role in Pipeline} & \textbf{Tradeoff Direction} \\
		\midrule
		Motion threshold & Separates background variation from candidate foreground & Higher values reduce noise but risk missed weak impacts \\
		Minimum blob area & Rejects tiny artifacts and sensor noise & Higher values improve precision but can remove valid small tips \\
		Morphology kernels & Connect fragmented dart regions before contouring & Larger kernels improve continuity but can over-merge nearby regions \\
		Still-time gate & Waits for stable post-impact frame before analysis & Longer gating reduces blur artifacts but increases response delay \\
		Known-dart suppression radius & Avoids repeated detection of already lodged darts & Larger radius reduces duplicates but may hide nearby new impacts \\
		\bottomrule
	\end{tabularx}
\end{center}

\vspace{\baselineskip}

\begin{lstlisting}[numbers=left,caption=Parameter tuning strategy used during detector stabilization, label=lst:parameter_tuning_strategy]
Start from conservative threshold and high min-blob area.
Verify no false triggers on static board for several seconds.
Lower threshold until weak impacts become visible in mask.
Re-tune morphology to reconnect fragmented dart regions.
Validate edge and out-of-bounds throws before locking values.
\end{lstlisting}

\setauthor{\firstauthor}
\section{Frontend and Visualization Technologies}
The frontend technology selection was guided by the need for rapid development, tight integration with reactive state management, and efficient runtime performance on widely distributed client machines. SvelteKit was selected because it compiles declarative component definitions into optimized vanilla JavaScript, avoiding the runtime bundle size overhead associated with virtual DOM frameworks such as React and Vue. The framework's first-class support for file-based routing, server-side rendering, and progressive enhancement enabled a concise codebase while preserving cross-platform behavior on Windows and Linux. Other frameworks, including Angular and Next.js, were considered but rejected due to their comparatively greater bundle complexity and configuration surface for a project with a single-page-game-like workflow.

Real-time communication is implemented using WebSocket transport, which provides low-latency, bidirectional messaging ideally suited for dart hit events generated by the detection backend. Alternative methods, such as polling over HTTP, Server-Sent Events (SSE), or WebRTC data channels, were evaluated. Polling was dismissed due to its increased latency and superfluous backend load, SSE was considered but does not support client-to-server messaging without an additional transport, and WebRTC exposed unnecessary complexity for this use case. The WebSocket approach balances implementation simplicity, responsiveness, and compatibility with the TypeScript ecosystem.

In the context of a desktop-only solution, Python GUI frameworks (e.g., PyQt, Tkinter) and Java UI toolkits (e.g., JavaFX, Swing) could have been viable. However, a browser-based architecture using SvelteKit offers broader compatibility, simpler deployment (no native installers), and access to modern web-oriented visualization libraries, which aligns with the requirement for cross-platform portability.

\subsection{Rendering}
The visualization stack relies on HTML canvas for custom game overlays and Chart.js for statistical plots. The canvas API provides deterministic pixel-level rendering for the virtual dartboard and enables efficient real-time drawing of dynamic elements such as hit markers and calibration guides. Chart.js was chosen for statistics visualization because it supports common chart types (bar, line, pie) with minimal configuration, interactive tooltips, and responsive design, which are adequate for the anticipated analytical requirements.

Alternative rendering technologies include SVG-based rendering (e.g., D3.js) and WebGL-based engines (e.g., PixiJS). SVG can be useful for vector-rich diagrams but suffers performance degradation with frequent updates, while WebGL provides high throughput at the cost of increased implementation complexity. For this project, the HTML canvas + Chart.js combination offers a pragmatic tradeoff between performance, ease of development, and maintainability.

\section{Backend and Data Management}

\subsection{Database Selection}
Database system selection was driven by the requirement that the application operate as a single-user local deployment with simple persistent storage semantics. SQLite meets these constraints by providing a lightweight, file-based relational database engine with ACID guarantees and zero administration overhead. Unlike client-server databases (e.g., PostgreSQL, MySQL), SQLite does not require standalone process management or inter-process communication. The expected data volume (game sessions, player records, score events) is modest, and SQLite performs well for this load while minimizing the risk of network partitioning or authentication misconfiguration in a local evaluation setting.

\subsection{Backend API Selection}
FastAPI was chosen as the backend framework due to its combination of asynchronous request handling, automatic OpenAPI schema generation, and modern Python typing support. In this system, the frontend is implemented with SvelteKit, which includes a native web server capability for static assets and route handling; therefore the backend API is strictly responsible for data and event services. FastAPI's lightweight routing and async I/O characteristics complement this architecture by enabling the backend to run as a dedicated API service without duplicating the web-serving responsibilities already managed by the SvelteKit process. Compared to Flask, FastAPI offers better scaling for concurrent I/O-bound workloads with built-in async/await support while maintaining similar engineering ergonomics. Compared to Django, FastAPI is more lightweight for microservice-style APIs and avoids the heavier template layer and ORM defaults that are unnecessary for the requirements of this project.

Regarding ORM selection, Tortoise ORM was adopted for its async-first design and integration with FastAPI. Alternative ORMs such as SQLAlchemy (classic) and Django ORM were evaluated; they provide rich capabilities but typically involve more boilerplate or blocking synchronous operations unless additional async adapters are used. Tortoise ORM offers a streamlined, declarative model syntax with async query execution, aligning with the project's asynchronous WebSocket and REST API architecture.

\setauthor{\secondauthor}
\section{Alternative Technology Approaches}

\subsection{Machine Learning-Based Detection}

\subsubsection{Data Requirements}
A learned detector would require a curated dataset of dart impacts across different lighting conditions, board wear states, dart geometries, and camera setups. Building such a dataset at the required scale was outside the scope of this project. Publicly available alternatives were either inaccessible or too narrow to support robust generalization, so a data-driven approach was not selected as the primary method.

\subsubsection{Generalization Challenges}
Even with sufficient data, model robustness would still need to be validated across sensor resolutions, focal lengths, and lens distortions, all of which change the apparent dart-tip shape. The chosen classical pipeline is less sensitive to appearance variation because it relies primarily on temporal change and geometric constraints rather than learned visual features.

\subsection{Multi-Camera Systems}

\subsubsection{Accuracy Improvements}
A two-camera setup could triangulate impact position in three dimensions and reduce ambiguity from single-view projection. It would also improve resilience to partial occlusion. However, this additional accuracy is not essential for the recreational scenario targeted in this thesis.

\subsubsection{Calibration Overhead}
Beyond per-camera intrinsic calibration, a stereo configuration requires stable extrinsic calibration between both cameras. Small mechanical shifts in mounts introduce baseline errors that directly degrade triangulation accuracy. This increases setup complexity and maintenance effort relative to a single-camera installation.

\section{Technology Selection Rationale}

\subsection{Justification for the Chosen Stack}

\subsubsection{Alignment with Project Objectives}
The selected stack was chosen to satisfy real-time performance, low-cost deployment, and maintainability with minimal hardware assumptions. OpenCV provides the required primitives for calibration, homography, thresholding, and contour analysis, while NumPy enables efficient numeric operations on shared array structures. Python provides rapid iteration and sufficient runtime performance because compute-intensive operations execute in compiled OpenCV routines. ArUco-based calibration supports repeatable setup in fixed installations, and the complete stack runs on Linux and Windows without specialized hardware beyond a standard camera.